\begingroup
\renewcommand{\clearpageforsection}{}
\exercisesheader{}
\endgroup

% 1

\eoce{\qt{Rentang hidup mamalia\label{mammal_life_spans}} Data rentang hidup (dalam 
tahun) dan masa kehamilan (dalam hari) dikumpulkan untuk 62 mamalia. Diagram pencar rentang 
hidup terhadap masa kehamilan ditampilkan di bawah ini. \footfullcite{Allison+Cicchetti:1975}

\noindent\begin{minipage}[c]{0.44\textwidth}
\begin{parts}
\item Jenis asosiasi apa yang tampak antara rentang hidup dan masa kehamilan?
\item Jenis asosiasi apa yang Anda perkirakan akan terlihat jika sumbu-sumbu diagram dibalik, yaitu jika masa kehamilan diplot terhadap rentang hidup?
\item Apakah rentang hidup dan masa kehamilan saling bebas? Jelaskan alasan Anda.
\end{parts}
\end{minipage}
\begin{minipage}[c]{0.55\textwidth}
\begin{center}
\Figures[Sebuah diagram pencar berisi 55 titik ditampilkan. Variabel "Masa kehamilan" ditampilkan pada sumbu horizontal dengan rentang dari 0 hari hingga sekitar 650 hari. Variabel "Rentang hidup" ditampilkan pada sumbu vertikal dengan rentang dari 0 tahun hingga 100 tahun. Sekumpulan besar titik berada antara 0 dan 250 hari kehamilan serta antara 0 dan 30 tahun. Di luar kumpulan ini, terdapat satu titik pada sekitar (10, 50). Kumpulan titik lain berada antara 250 dan 450 hari kehamilan serta antara 25 dan 50 tahun. Di luar titik-titik yang telah dijelaskan, terdapat tiga titik pada (250 hari, 100 tahun), (640 hari, 70 tahun), dan (650 hari, 45 tahun).]{0.86}{eoce/mammal_life_spans}{mammal_life_spans_scatterplot}
\end{center}
\end{minipage}
}{}

% 2

\eoce{\qt{Asosiasi\label{association_plots}}
Tentukan diagram mana yang menunjukkan
(a)~asosiasi positif,
(b)~asosiasi negatif, atau
(c)~tidak ada asosiasi.
Tentukan juga apakah asosiasi positif dan negatif tersebut
linear atau nonlinear.
Setiap bagian dapat merujuk pada lebih dari satu diagram.
\begin{center}
\Figures[Empat diagram pencar ditampilkan dan diberi label 1, 2, 3, dan 4. Sumbu-sumbunya tidak berlabel karena hanya pola titik pada diagram yang penting untuk latihan ini. Pada diagram 1, titik-titik cukup mengelompok di sudut kiri bawah dan tetap mengelompok ketika bergerak ke kanan, dengan pola yang naik secara teratur menuju sudut kanan atas. Pada diagram 2, titik-titik tampak tersebar hampir secara acak di seluruh daerah gambar berbentuk persegi panjang. Diagram 3 menunjukkan titik-titik yang mengelompok rapat di sudut kiri bawah dan tetap mengelompok ketika bergerak ke kanan, dengan pola yang mula-mula naik secara bertahap lalu semakin curam ketika mencapai sisi kanan diagram. Pada bagian kiri diagram 4, data cukup mengelompok di sudut kiri atas, lalu menurun secara teratur menuju sudut kanan bawah.]{0.95}{eoce/association_plots}{association_plots}
\end{center}
}{}

% 3

\eoce{\qt{Bakteri berkembang biak\label{reproducing_bacteria}} Andaikan ruang dan nutrisi 
dalam sebuah cawan petri hanya cukup untuk menopang satu juta sel bakteri. Anda menempatkan 
beberapa sel bakteri di dalam cawan petri tersebut, membiarkannya berkembang biak tanpa 
hambatan, dan mencatat jumlah sel bakteri di dalam cawan dari waktu ke waktu. Buatlah sketsa 
diagram yang menggambarkan hubungan antara jumlah sel bakteri dan waktu.
% first exponential
}{}

% 4

\eoce{\qt{Produktivitas kantor\label{office_productivity}} Produktivitas kantor relatif rendah 
ketika karyawan sama sekali tidak merasa tertekan oleh pekerjaan atau keamanan kerja mereka. 
Namun, tingkat stres yang tinggi juga dapat menurunkan produktivitas karyawan. Buatlah sketsa 
diagram yang menggambarkan hubungan antara stres dan produktivitas.
}{}

% 5

\eoce{\qt{Parameter dan statistik\label{parameters_stats}} Tentukan nilai mana yang merupakan 
rata-rata sampel dan nilai mana yang merupakan rata-rata populasi yang diklaim.
\begin{parts}
\item Pada tahun 2007, rumah tangga di Amerika Serikat membelanjakan rata-rata sekitar \$52 
untuk keperluan Halloween seperti kostum, dekorasi, dan permen. Untuk mengetahui apakah 
angka ini telah berubah, peneliti melakukan survei baru pada tahun 2008 sebelum angka dari 
industri dilaporkan. Survei tersebut mencakup 1,500 rumah tangga dan menemukan bahwa 
pengeluaran Halloween rata-rata adalah \$58 per rumah tangga.
\item Rata-rata indeks prestasi kumulatif (GPA) mahasiswa pada tahun 2001 di sebuah universitas 
swasta adalah 3.37. Satu dekade kemudian, survei terhadap sampel yang terdiri atas 203 mahasiswa 
dari universitas tersebut menghasilkan GPA rata-rata sebesar 3.59.
\end{parts}
}{}

% 6

\eoce{\qt{Tidur semasa kuliah\label{college_sleeping}} Sebuah artikel terbaru di surat kabar 
kampus menyatakan bahwa mahasiswa tidur rata-rata 5.5 jam setiap malam. Seorang mahasiswa 
yang meragukan nilai ini memutuskan untuk melakukan survei dengan mengambil sampel acak 
sebanyak 25 mahasiswa. Mahasiswa dalam sampel tersebut tidur rata-rata 6.25 jam per malam. 
Tentukan nilai mana yang merupakan rata-rata sampel dan nilai mana yang merupakan rata-rata 
populasi yang diklaim.
}{}

% R011-B005 D021: jeda halaman paksa untuk cetak ditiadakan.

% 7

\eoce{\qt{Hari libur di lokasi pertambangan\label{days_off_mining}} Para pekerja di sebuah 
lokasi pertambangan menerima rata-rata 35 hari cuti berbayar, lebih rendah daripada rata-rata 
nasional. Manajer lokasi tersebut mendapat tekanan dari serikat pekerja setempat untuk menambah 
jumlah hari cuti berbayar. Namun, ia tidak ingin memberi pekerja lebih banyak hari libur karena 
hal itu akan mahal. Sebagai gantinya, ia memutuskan untuk memecat 10 karyawan sedemikian rupa 
sehingga rata-rata jumlah hari libur yang dilaporkan oleh para karyawannya meningkat. Untuk 
mencapai tujuan ini, apakah ia harus memecat karyawan yang memiliki hari libur terbanyak, 
yang memiliki hari libur paling sedikit, atau yang memiliki jumlah hari libur sekitar rata-rata?
}{}

% 8

\eoce{\qt{Median dan IQR} Untuk setiap bagian, bandingkan distribusi (1) dan (2) berdasarkan median dan IQR-nya. Anda tidak perlu menghitung statistik tersebut; cukup nyatakan bagaimana median dan IQR kedua distribusi dibandingkan. Pastikan untuk menjelaskan alasan Anda. 
\begin{multicols}{2}
\begin{parts}
\item (1) 3, 5, 6, 7, 9 \\
(2) 3, 5, 6, 7, 20
\item (1) 3, 5, 6, 7, 9 \\
(2) 3, 5, 7, 8, 9
\item (1) 1, 2, 3, 4, 5 \\
(2) 6, 7, 8, 9, 10
\item (1) 0, 10, 50, 60, 100 \\
(2) 0, 100, 500, 600, 1000
\end{parts}
\end{multicols}
}{}

% 9

\eoce{\qt{Rata-rata dan simpangan baku} Untuk setiap bagian, bandingkan distribusi (1) dan (2) berdasarkan rata-rata dan simpangan bakunya. Anda tidak perlu menghitung statistik tersebut; cukup nyatakan bagaimana rata-rata dan simpangan baku kedua distribusi dibandingkan. Pastikan untuk menjelaskan alasan Anda. \textit{Petunjuk:} Membuat sketsa diagram titik kedua distribusi mungkin berguna.
\begin{multicols}{2}
\begin{parts}
\item (1) 3, 5, 5, 5, 8, 11, 11, 11, 13 \\
(2) 3, 5, 5, 5, 8, 11, 11, 11, 20 \\
\item (1) -20, 0, 0, 0, 15, 25, 30, 30 \\
(2) -40, 0, 0, 0, 15, 25, 30, 30
\item (1) 0, 2, 4, 6, 8, 10 \\
(2) 20, 22, 24, 26, 28, 30
\item (1) 100, 200, 300, 400, 500 \\
(2) 0, 50, 300, 550, 600
\end{parts}
\end{multicols}
}{}

% 10

\eoce{\qt{Cocokkan} Deskripsikan distribusi pada histogram di bawah ini dan cocokkan setiap histogram dengan diagram kotaknya. \\
\begin{center}
\Figures[Enam diagram ditampilkan: tiga histogram berlabel a, b, dan c serta 3 diagram kotak berlabel 1, 2, dan 3.

Diagram (a) menunjukkan histogram dengan rentang data horizontal dari 50 hingga 70. Data berbentuk lonceng dan berpusat di tengah diagram; hanya sedikit data yang mendekati batas bawah 50 dan batas atas 70.

Diagram (b) menunjukkan histogram lain dengan sumbu horizontal yang membentang dari 0 hingga 100. Tinggi kelas-kelas interval histogram relatif tetap, mulai dari kelas pertama di dekat nol hingga kelas terakhir di dekat 100.

Diagram (c) adalah histogram dengan sumbu horizontal dari 0 hingga sekitar 7. Beberapa kelas interval pertama naik dengan cepat menuju puncak pada nilai horizontal 1, lalu turun hingga nilai 2. Setelah itu, tingginya menurun jauh lebih bertahap hingga sekitar 4, tempat kelas-kelas interval mendekati nol dan tetap mendekati nol untuk nilai yang lebih besar.

Diagram (1) adalah diagram kotak dengan sumbu vertikal dari 0 hingga sekitar 7. Garis kumis bawah berada pada 0, kotaknya membentang dari sekitar 1 hingga 2, dan garis tengah diagram kotak berada pada sekitar 1.4. Garis kumis atas mencapai sekitar 3.5, lalu beberapa titik ditandai secara terpisah hingga sekitar 7.

Diagram (2) adalah diagram kotak dengan sumbu vertikal yang membentang dari sekitar 50 hingga 70. Kotaknya berpusat pada 60 dan membentang dari sekitar 58 hingga 62. Garis kumis membentang dari sekitar 52 hingga 68. Sebanyak 2 titik terpisah ditampilkan di bawah 52 dan sekitar 4 titik ditampilkan di atas 68.

Diagram (3) adalah diagram kotak yang membentang dari 0 hingga 100. Kotaknya berpusat pada sekitar 50 dan membentang dari sekitar 25 hingga 75. Garis kumis menjulur ke bawah hingga 0 dan ke atas hingga 100.]{}{eoce/hist_box_match}{hist_box_match}
\end{center}
}{}

% R011-B005 D028: jeda halaman paksa setelah Latihan 2.10 ditiadakan
% agar Latihan 2.11 mengikuti panel gambar pada halaman yang sama.

% 11

\eoce{\qt{Kualitas udara\label{air_quality_durham}} Kualitas udara harian diukur dengan indeks 
kualitas udara (AQI) yang dilaporkan oleh Badan Perlindungan Lingkungan Amerika Serikat. Indeks 
ini melaporkan tingkat pencemaran dan dampak kesehatan terkait yang mungkin perlu diperhatikan. 
Indeks tersebut dihitung untuk lima pencemar udara utama yang diatur oleh Undang-Undang Udara 
Bersih dan memiliki nilai dari 0 hingga 300; nilai yang lebih tinggi menunjukkan kualitas udara 
yang lebih rendah. AQI dilaporkan untuk sampel sebanyak 91 hari pada tahun 2011 di Durham, 
Carolina Utara. Histogram frekuensi relatif di bawah ini menunjukkan distribusi nilai AQI pada 
hari-hari tersebut. \footfullcite{data:durhamAQI:2011} \\
\begin{minipage}[c]{0.55\textwidth}
\begin{parts}
\item Perkirakan nilai median AQI sampel ini.
\item Apakah Anda memperkirakan nilai rata-rata AQI sampel ini lebih tinggi atau lebih rendah 
daripada mediannya? Jelaskan alasan Anda.
\item Perkirakan Q1, Q3, dan IQR distribusi tersebut.
\item Apakah ada hari dalam sampel ini yang dapat dianggap memiliki AQI yang sangat rendah atau 
sangat tinggi? Jelaskan alasan Anda.
\end{parts}
\end{minipage}
\begin{minipage}[c]{0.45\textwidth}
\begin{center}
\Figures[Sebuah histogram "AQI harian" ditampilkan dengan sumbu horizontal data dari sekitar 5 hingga 65. Lebar kelas intervalnya 5, terdapat 12 kelas interval dari 5 hingga 60, dan sumbu vertikal menunjukkan proporsi. Tinggi ke-12 kelas interval tersebut, secara berurutan dari kiri ke kanan, sekitar 0.02 (untuk kelas 5 hingga 10), 0.06, 0.20, 0.06, 0.20, 0.15, 0.07, 0.04, 0.07, 0.08, 0.03, dan 0.02 untuk kelas terakhir dari 60 hingga 65.]{}{eoce/air_quality_durham}{air_quality_durham_rel_freq_hist}
\end{center}
\end{minipage}
}{}

% 12

\eoce{\qt{Median versus rata-rata\label{estimate_mean_median_simple}} Perkirakan median untuk 
400 pengamatan yang ditampilkan dalam histogram, lalu nyatakan apakah Anda memperkirakan 
rata-ratanya lebih tinggi atau lebih rendah daripada median.
\begin{center}
\Figures[Sebuah histogram ditampilkan dengan sumbu horizontal data dari 40 hingga 100 dan lebar kelas interval 5. Frekuensi kelas-kelas interval tersebut adalah sebagai berikut, dengan jumlah yang bersifat perkiraan: 2 (untuk kelas 40 hingga 45), 4, 2, 10, 20, 25, 50, 75, 70, 85, 45, dan 10 untuk kelas terakhir dari 95 hingga 100.
]{0.6}{eoce/estimate_mean_median_simple}{estimate_mean_median_simple}
\end{center}
}{}

% 13

\eoce{\qt{Histogram versus diagram kotak\label{hist_vs_box}} Bandingkan kedua diagram di bawah 
ini. Karakteristik distribusi apa yang tampak dalam histogram tetapi tidak dalam diagram kotak? 
Karakteristik apa yang tampak dalam diagram kotak tetapi tidak dalam histogram?
\begin{center}
\Figures[Dua diagram ditampilkan: yang pertama histogram dan yang kedua diagram kotak. Data pada setiap diagram membentang dari sekitar 0 hingga 30.

Histogram memiliki kelas interval dengan lebar 2. Dimulai dari nilai yang lebih rendah, kelas-kelas interval menunjukkan puncak pertama pada nilai horizontal sekitar 5, lalu menurun hingga mendekati sumbu horizontal pada 10, sebelum naik kembali antara 10 dan 14. Tinggi kelas-kelas interval kembali lebih rendah antara 15 dan 30.

Kotak pada diagram kotak berpusat pada 10 dan membentang dari sekitar 5 hingga 12. Garis kumisnya mencapai sekitar 2 di bawah dan sekitar 22 di atas. Terdapat beberapa titik di atas garis kumis atas.
]{0.6}{eoce/hist_vs_box}{hist_vs_box}
\end{center}
}{}

% 14

\eoce{\qt{Teman Facebook\label{dist_shape_fb_friends}} Data Facebook menunjukkan bahwa 
50\% pengguna Facebook memiliki 100 teman atau lebih dan bahwa jumlah teman rata-rata 
pengguna adalah 190. Apa yang disiratkan oleh temuan ini tentang bentuk distribusi jumlah 
teman pengguna Facebook? \footfullcite{Backstrom:2011}
}{}

% 15

\eoce{\qt{Distribusi dan statistik yang sesuai, Bagian I\label{dist_shape_pets_dist_height}} 
Untuk setiap hal berikut, nyatakan apakah Anda memperkirakan distribusinya simetris, menceng 
kanan, atau menceng kiri. Tentukan juga apakah rata-rata atau median paling tepat mewakili 
pengamatan tipikal dalam data, serta apakah variabilitas pengamatan paling tepat diwakili oleh 
simpangan baku atau IQR. Jelaskan alasan Anda.
\begin{parts}
\item Jumlah hewan peliharaan per rumah tangga. 
\item Jarak ke tempat kerja, yaitu jarak dalam mil antara tempat kerja dan rumah.
\item Tinggi badan laki-laki dewasa.
\end{parts}
}{}

\D{\newpage}

% 16

\eoce{\qt{Distribusi dan statistik yang sesuai, Bagian II\label{dist_shape_housing_alcohol_salary}} 
Untuk setiap hal berikut, nyatakan apakah Anda memperkirakan distribusinya simetris, menceng 
kanan, atau menceng kiri. Tentukan juga apakah rata-rata atau median paling tepat mewakili 
pengamatan tipikal dalam data, serta apakah variabilitas pengamatan paling tepat diwakili oleh 
simpangan baku atau IQR. Jelaskan alasan Anda.
\begin{parts}
\item Harga rumah di sebuah negara tempat 25\% rumah berharga di bawah \$350,000, 
50\% rumah berharga di bawah \$450,000, 75\% rumah berharga di bawah \$1,000,000, 
dan cukup banyak rumah berharga lebih dari \$6,000,000.
\item Harga rumah di sebuah negara tempat 25\% rumah berharga di bawah \$300,000, 
50\% rumah berharga di bawah \$600,000, 75\% rumah berharga di bawah \$900,000, 
dan sangat sedikit rumah berharga lebih dari \$1,200,000.
\item Jumlah minuman beralkohol yang dikonsumsi mahasiswa dalam satu minggu. Andaikan 
sebagian besar mahasiswa tersebut tidak minum karena berusia di bawah 21 tahun dan hanya 
sedikit yang minum secara berlebihan.
\item Gaji tahunan karyawan di sebuah perusahaan Fortune 500, tempat hanya beberapa 
eksekutif tingkat tinggi yang memperoleh gaji jauh lebih besar daripada semua karyawan lain.
\end{parts}
}{}

% 17

\eoce{\qt{Pendapatan di kedai kopi\label{income_coffee_shop}} Histogram pertama di bawah ini 
menunjukkan distribusi pendapatan tahunan 40 pelanggan di sebuah kedai kopi kampus. Andaikan 
dua orang baru memasuki kedai kopi tersebut: seorang berpenghasilan \$225,000 dan seorang lagi 
berpenghasilan \$250,000. Histogram kedua menunjukkan distribusi pendapatan yang baru. Statistik 
ringkasan juga disediakan. \\
\begin{minipage}[c]{0.57\textwidth}
\Figures[Dua histogram ditampilkan dan diberi label 1 dan 2. Diagram 1 memiliki sumbu horizontal dari \$60,000 hingga \$70,000. Dari kiri ke kanan, tinggi kelas-kelas interval umumnya meningkat secara teratur dari frekuensi 2 hingga 3 pada \$60,000 sampai \$62,000, lalu mencapai puncak sekitar 7 hingga 8 antara \$64,000 dan \$66,000. Setelah itu, jumlah pada kelas-kelas interval menurun secara teratur hingga sekitar 2 pada kelas terakhir, yaitu \$69,000 sampai \$70,000. Diagram (2) menunjukkan histogram dengan sumbu horizontal dari sekitar \$60,000 hingga \$260,000. Lebar kelas intervalnya \$1,000 seperti pada diagram pertama, dan 10 kelas interval pertama menggambarkan data yang dijelaskan pada Diagram (1). Dua kelas interval tambahan ditampilkan pada sekitar \$225,000 dan \$250,000, masing-masing dengan tinggi 1.]{}{eoce/income_coffee_shop}{income_coffee_shop}
\end{minipage}
\begin{minipage}[c]{0.4\textwidth}
\begin{center}
\begin{tabular}{rrr}
\hline
             & (1)       & (2) \\ 
\hline
n            & 40        & 42 \\ 
Min.         & 60,680    & 60,680 \\ 
Ku. 1        & 63,620    & 63,710 \\ 
Median       & 65,240    & 65,350 \\ 
Rata-rata    & 65,090    & 73,300 \\ 
Ku. 3        & 66,160    & 66,540 \\ 
Maks.        & 69,890    & 250,000 \\ 
Simp. baku   & 2,122     & 37,321 \\ 
\hline
\end{tabular}
\end{center}
\end{minipage}
\begin{parts}
\item Apakah rata-rata atau median paling tepat mewakili pendapatan yang dapat kita anggap 
tipikal bagi 42 pelanggan kedai kopi ini? Apa yang ditunjukkan hal tersebut tentang sifat 
robust kedua ukuran ini?
\item Apakah simpangan baku atau IQR paling tepat mewakili besarnya variabilitas pendapatan 
42 pelanggan kedai kopi ini? Apa yang ditunjukkan hal tersebut tentang sifat robust kedua ukuran ini?
\end{parts}
}{}

% 18

\eoce{\qt{Nilai tengah rentang\label{midrange}} \textit{Nilai tengah rentang} suatu distribusi 
didefinisikan sebagai rata-rata dari nilai maksimum dan minimum distribusi tersebut. Apakah 
statistik ini robust terhadap pencilan dan kemencengan ekstrem? Jelaskan alasan Anda.
}{}

\D{\newpage}

% 19

\eoce{\qt{Waktu perjalanan ke tempat kerja\label{county_commute_times}} Sensus Amerika Serikat 
mengumpulkan data tentang waktu yang dibutuhkan penduduk Amerika untuk bepergian ke tempat 
kerja, selain banyak variabel lain. Histogram di bawah ini menunjukkan distribusi waktu 
perjalanan rata-rata di 3,142 county di Amerika Serikat pada tahun 2010. Peta intensitas 
spasial dari data yang sama juga ditampilkan di bawah ini.
\begin{center}
\Figures[Sebuah histogram ditampilkan dengan sumbu horizontal untuk variabel "Rerata waktu perjalanan kerja (menit)" yang membentang dari sekitar 0 hingga 50. Sumbu vertikal menunjukkan frekuensi dengan nilai puncak sekitar 200. Kelas-kelas interval pada sisi kiri mula-mula rendah, lalu tingginya mulai meningkat pada sekitar 10 dan naik dengan cepat menjelang 15 sebelum melandai dan mencapai puncak pada sekitar 22. Tinggi kelas-kelas interval mulai menurun lagi secara bertahap pada sekitar 24, kemudian lebih cepat pada sekitar 26 hingga 29. Pada 30, tingginya terus menurun tetapi dengan laju yang lebih lambat, lalu mendatar mendekati 0 pada sekitar 35.]{0.48}{eoce/county_commute_times}{county_commute_times_hist}
\Figures[Sebuah peta intensitas spasial Amerika Serikat ditampilkan. Legenda arsiran membentang dari nilai 4 hingga sekitar 33. Arsiran di bagian timur negara tersebut menunjukkan nilai yang sedikit lebih tinggi, sedangkan bagian barat dari Midwest bagian utara (Dakota Utara, Dakota Selatan, dan Nebraska) menunjukkan nilai yang lebih rendah. Wilayah khusus lain yang menunjukkan pola nilai lebih tinggi daripada daerah sekitarnya berada di Florida bagian selatan dan California bagian utara.]{0.48}{eoce/county_commute_times}{county_commute_times_map}
\end{center}
\begin{parts}
\item Deskripsikan distribusi numeriknya dan berikan pendapat apakah transformasi logaritma 
layak diterapkan pada data ini.
\item Deskripsikan distribusi spasial waktu perjalanan ke tempat kerja dengan menggunakan peta di atas.
\end{parts} 
}{}

% 20

\eoce{\qt{Populasi Hispanik\label{county_hispanic_pop}} Sensus Amerika Serikat mengumpulkan 
data tentang ras dan etnis penduduk Amerika, selain banyak variabel lain. Histogram di bawah 
ini menunjukkan distribusi persentase penduduk Hispanik di 3,142 county di Amerika Serikat 
pada tahun 2010. Histogram logaritma nilai-nilai tersebut juga ditampilkan.
\begin{center}
\Figures[Sebuah histogram untuk variabel "Persentase Hispanik" ditampilkan dengan sumbu horizontal dari 0 hingga 100. Kelas interval pertama, dari 0 hingga 5, jauh lebih tinggi daripada semua kelas lain, yaitu sekitar 2000. Setelah itu, tinggi kelas-kelas interval turun dengan cepat: sekitar 500 antara 5 dan 10, 200 antara 10 dan 15, serta 100 antara 15 dan 20. Selanjutnya, tingginya terus menurun hingga hampir tidak dapat dibedakan dari 0 untuk kelas-kelas interval di sekitar 50\%.]{0.48}{eoce/county_hispanic_pop}{county_hispanic_pop_hist}
\Figures[Sebuah histogram untuk variabel yang telah ditransformasi, "logaritma natural Persentase Hispanik", ditampilkan dengan sumbu horizontal dari sekitar -2.5 hingga 4.5. Frekuensi kelas-kelas interval sangat dekat dengan 0 hingga -1, lalu naik sedikit sampai sekitar -0.5 sebelum meningkat tajam menuju puncak pada sekitar 0.5. Setelah itu, tingginya menurun secara teratur menuju frekuensi 0 pada nilai horizontal 4.5.]{0.48}{eoce/county_hispanic_pop}{county_hispanic_pop_log_hist}
\Figures[Sebuah peta intensitas spasial Amerika Serikat ditampilkan. Legenda arsiran membentang dari nilai 0\% hingga nilai tertinggi "lebih dari 40\%". Sebagian besar wilayah timur dan tengah negara tersebut---di sebelah timur Texas, Colorado, Utah, dan Idaho---diarsir dengan nilai di bawah 10\%. Florida merupakan pengecualian karena beberapa county menunjukkan nilai yang lebih tinggi. Nilai yang lebih tinggi sangat menonjol di Texas, New Mexico, Arizona, dan California, yang sebagian besar menunjukkan arsiran untuk nilai sedikitnya 20\%. Nevada, Idaho, Oregon, dan Washington menunjukkan nilai rata-rata sekitar 10--20\%.]{0.48}{eoce/county_hispanic_pop}{county_hispanic_pop_map}
\end{center}
\begin{parts}
\item Deskripsikan distribusi numeriknya dan jelaskan mengapa kita mungkin ingin menggunakan 
nilai yang telah ditransformasi dengan logaritma ketika menganalisis atau memodelkan data ini.
\item Karakteristik distribusi populasi Hispanik di county-county Amerika Serikat apa yang 
tampak pada peta tetapi tidak pada histogram? Karakteristik apa yang tampak pada histogram 
tetapi tidak pada peta?
\item Apakah salah satu visualisasi lebih tepat atau lebih membantu daripada yang lain? 
Jelaskan alasan Anda.
\end{parts} 
}{}
